\section{Example of an Appendix}
\addcontentsline{toc}{section}{Anexo A: Exa,ple of Appendix}
Appendices may be used to provide supplementary material that is relevant to the paper but would interrupt the flow of the main text if included in the body of the manuscript. Typical examples include detailed derivations, additional tables, implementation notes, or complementary experimental results.

%As an illustration, suppose that an intermediate result used in the main text is given by
%\begin{equation}
%y[n] = \sum_{k=0}^{M-1} h[k]x[n-k].
%\label{eq_appendix_conv}
%\end{equation}
%Equation~\eqref{eq_appendix_conv} represents a discrete-time convolution and is included here only as a formatting example.
%
%Authors should include appendices only when they add useful information for the reader. If no appendix is needed, this part of the template may simply be removed. If the paper is an Extended Abstract, appendices are not allowed.
